\subsection{Sphericality and Galois reduction at multiplicity three}
\label{sec:frontier-spherical-reduction}
\begin{proposition}\label{prop:frontier-galois-pu}
Every complex fusion category with based ring $\RR{3}{215}$ admits a
spherical structure and has a pseudo-unitary Galois conjugate.
Consequently general categorifiability of this ring is equivalent to
pseudo-unitary categorifiability. This reduction alone asserts neither
existence nor nonexistence; the final obstruction is proved in
Theorem~\ref{thm:frontier-noncategorifiable}.
\end{proposition}
\begin{proof}
Apply the exhaustive oriented signed-sphericalization construction of
Section~\ref{sec:signed}, using every coefficient in its full parity
interval and imposing no positive-indicator or modulo-four normalization.
There are $1480$ associative signed quotients. The $32$ unit-preserving
sign gauges partition them into $63$ orbits; one orbit consists of the
$32$ gauges of the original fusion tensor. On representatives of the
other $62$ orbits there are $271$ rational character components.
Of these, $228$ have a zero simple dimension and are impossible.
For each of the other $43$, the norm of the prospective global
dimension has $17$-adic valuation zero. The original ring has a formal
codegree with minimal polynomial $t^2-272t+1088$, whose norm has
$17$-adic valuation one. Thus the signed-sphericalization
norm-divisibility condition rejects every remaining component. No
nongauge quotient can occur for an actual category.

For completeness, a gauge quotient gives sphericality of the original
category, not merely compatibility with spherical dimensions. If
$T_{ij}^k=\eta_i\eta_j\eta_kN_{ij}^k$, replace the chosen lift of $X_i$
in the sphericalization by its tensor product with the sign object
when $\eta_i=-1$. The new signed quotient has $T=N$. The sum and
difference of the multiplicities of the two lifts of $X_k$ in a tensor
product are then both $N_{ij}^k$, so the opposite lift has multiplicity
zero. These chosen lifts therefore generate a full tensor subcategory
of the sphericalization which projects tensor-equivalently onto the
original category. Transport its spherical structure along that
equivalence.

The four rational characters of the original ring have a zero entry,
so none is a spherical dimension character. The only remaining
character orbit is
\[
 (1,1,d,d,d-1,d+1),\qquad d^2-8d-1=0.
\]
One of its two real embeddings is the Frobenius--Perron character,
$d=4+\sqrt{17}$. Galois conjugation to that embedding gives a spherical
category with positive Frobenius--Perron dimensions, hence a
pseudo-unitary category. Conversely a pseudo-unitary categorification
is a complex fusion categorification.

The complete unrestricted signed enumeration has $274655$ search nodes
and is frozen in
\filename{research/new_frontier/frontier_spherical_reduction.json}.
The companion Python script recomputes all $1480$ quotients and the
complete sign orbits, verifies the original gauge orbit, checks all
$271$ nongauge character components and their arithmetic obstructions,
and reconstructs the original dimension-character orbit using exact
arithmetic. The older normalized twelve-table enumeration is retained
separately as a regression control; it is not needed for this proof.

\end{proof}
